\documentclass[11pt,a4paper]{article}
\usepackage[T1]{fontenc}
\usepackage[utf8]{inputenc}
\usepackage[french]{babel}
\usepackage{lmodern}
\usepackage{textcomp}
\usepackage{amsmath,amssymb}
\usepackage{geometry}
\usepackage{tabularx,array}
\usepackage[table]{xcolor}
\usepackage{tikz}
\geometry{left=8mm,right=8mm,top=7mm,bottom=7mm}
\setlength{\parindent}{0pt}
\pagestyle{empty}
\setlength{\tabcolsep}{3pt}
\renewcommand{\arraystretch}{1.05}
\renewcommand{\tabularxcolumn}[1]{m{#1}}

% --- Palette pensée pour une impression noir et blanc ---
\definecolor{canlight}{gray}{0.93}
\definecolor{canmid}{gray}{0.72}
\definecolor{candark}{gray}{0.22}
\definecolor{answerblue}{RGB}{0,82,165}

% --- Colonnes ---
\newcolumntype{Q}{>{\centering\arraybackslash}X}
\newcolumntype{N}{>{\centering\arraybackslash}m{6mm}}
\newcolumntype{R}{>{\centering\arraybackslash}m{22mm}}

% --- En-tête ---
\newcommand{\canheader}[1]{%
  \begin{center}
    {\LARGE\bfseries COURSE AUX NOMBRES}\par
    \vspace{0.7mm}
    {\color{canmid}\rule{27mm}{0.45pt}\hspace{2mm}\color{black}$\diamond$\hspace{1.4mm}{\large\bfseries #1}\hspace{1.4mm}$\diamond$\hspace{2mm}\color{canmid}\rule{27mm}{0.45pt}\color{black}}
  \end{center}
  \vspace{0.8mm}
  \renewcommand{\arraystretch}{1.25}%
  \begin{tabularx}{\linewidth}{>{\centering\arraybackslash}X >{\centering\arraybackslash}X >{\centering\arraybackslash}m{40mm}}
    \rowcolor{canlight}
    \footnotesize\textbf{30 questions} -- 9 minutes &
    \footnotesize\textbf{Sans calculatrice} &
    \footnotesize\textbf{Score :} \rule{13mm}{0.45pt}\,/30
  \end{tabularx}
  \renewcommand{\arraystretch}{1.05}%
  \vspace{1mm}
  \begin{center}
    {\scriptsize Réponse uniquement \quad$\diamond$\quad Reporter le score sur la fiche annuelle.}
  \end{center}
  \vspace{1.2mm}
}

\newcommand{\qrow}[2]{\textbf{#1} & \parbox[c][15.0mm][c]{\linewidth}{\centering #2} & \\ \hline}
\newcommand{\qrowfig}[2]{\textbf{#1} & \parbox[c][27mm][c]{\linewidth}{\centering #2} & \\ \hline}
\newcommand{\qrowcorr}[3]{\textbf{#1} & \parbox[c][15.0mm][c]{\linewidth}{\centering #2} & {\color{answerblue}\bfseries #3} \\ \hline}
\newcommand{\qrowfigcorr}[3]{\textbf{#1} & \parbox[c][27mm][c]{\linewidth}{\centering #2} & {\color{answerblue}\bfseries #3} \\ \hline}

\newcommand{\cantable}[1]{%
\arrayrulecolor{canmid}
\begin{tabularx}{\linewidth}{N|Q|R}
\rowcolor{canlight}\textbf{N\textsuperscript{o}}&\centering\textbf{Question}&\textbf{Réponse}\\\hline
#1
\end{tabularx}
\arrayrulecolor{black}
}

\newcommand{\hypotfigure}{%
\begin{tikzpicture}[baseline=(current bounding box.center),scale=0.72,line width=0.6pt]
  \node[font=\bfseries\normalsize] at (1.6,2.55) {Quelle est l'hypoténuse ?};
  \coordinate (A) at (0,0); \coordinate (B) at (0,1.85); \coordinate (C) at (3.2,0);
  \draw (A)--(B)--(C)--cycle;
  \draw (0,0) rectangle (0.32,0.32);
  \node[below left] at (A) {\small A};
  \node[above left] at (B) {\small B};
  \node[below right] at (C) {\small C};
\end{tikzpicture}}


\newcommand{\canheadercorr}[1]{%
  \begin{center}
    {\LARGE\bfseries COURSE AUX NOMBRES}\par
    \vspace{0.7mm}
    {\color{canmid}\rule{27mm}{0.45pt}\hspace{2mm}\color{black}$\diamond$\hspace{1.4mm}{\large\bfseries #1}\hspace{1.4mm}$\diamond$\hspace{2mm}\color{canmid}\rule{27mm}{0.45pt}\color{black}}
  \end{center}
  \vspace{0.8mm}
  \renewcommand{\arraystretch}{1.25}%
  \begin{tabularx}{\linewidth}{>{\centering\arraybackslash}X >{\centering\arraybackslash}X >{\centering\arraybackslash}m{40mm}}
    \rowcolor{canlight}
    \footnotesize\textbf{30 questions} -- 9 minutes &
    \footnotesize\textbf{Sans calculatrice} &
    \footnotesize\textcolor{answerblue}{\textbf{CORRIGÉ}}
  \end{tabularx}
  \renewcommand{\arraystretch}{1.05}%
  \vspace{1mm}
  \begin{center}
    {\scriptsize Les réponses apparaissent en \textcolor{answerblue}{\textbf{bleu}}.}
  \end{center}
  \vspace{1.2mm}
}


\begin{document}
\canheader{CAN 2}

\begin{minipage}[t]{0.493\textwidth}
\normalsize
\cantable{
\qrow{1}{$72\div8$}
\qrow{2}{Écrire en chiffres : quinze unités et trois centièmes.}
\qrow{3}{$5-(-2)$}
\qrow{4}{Compléter : $47+\ldots=100$}
\qrow{5}{$10\,\%$ de 350}
\qrow{6}{La différence entre le double de 9 et le tiers de 12}
\qrow{7}{$0{,}8+1{,}35$}
\qrow{8}{Quel est le nombre qui, multiplié par 3, donne 7 ? Donner une fraction.}
\qrow{9}{$3{,}6\div100$}
\qrow{10}{Dans la série 4 ; 6 ; 6 ; 7 ; 9, quel est l'effectif de 6 ?}
\qrow{11}{Le plus grand : $-3$ ou $-7$ ?}
\qrow{12}{Quelle est la distance à zéro de $-6$ ?}
\qrow{13}{$\dfrac25$ de 35}
\qrow{14}{135 min = $\ldots$ h $\ldots$ min}
\qrow{15}{$(18-6)\div3$}
}
\end{minipage}\hfill
\begin{minipage}[t]{0.493\textwidth}
\normalsize
\cantable{
\qrow{16}{Le plus grand : $\dfrac34$ ou $\dfrac58$ ?}
\qrow{17}{Aire d'un rectangle de 8 cm sur 3 cm}
\qrow{18}{Moyenne de 7 ; 8 ; 9 ; 12}
\qrow{19}{Compléter : $1{,}2=\dfrac{\ldots}{10}$}
\qrow{20}{$1\,\mathrm{m}^2$ correspond à combien de carrés de $1\,\mathrm{cm}^2$ ?}
\qrow{21}{Quel est le reste de la division de 29 par 4 ?}
\qrow{22}{$2{,}5\times4$}
\qrow{23}{$4\,500$ cm = $\ldots$ m}
\qrow{24}{4 cahiers coûtent 6 \texteuro. Combien coûtent 10 cahiers ?}
\qrow{25}{Dans un triangle, deux angles mesurent $90^\circ$ et $35^\circ$. Le troisième mesure :}
\qrow{26}{$11^2$}
\qrow{27}{Un carré de côté 8 cm peut-il avoir une aire de $70\,\mathrm{cm}^2$ ?}
\qrow{28}{Réduire : $x+x+x$}
\qrow{29}{$5(x+2)$ est-il une somme ou un produit ?}
\qrow{30}{La température passe de $-4\,^{\circ}\mathrm C$ à $3\,^{\circ}\mathrm C$. De combien augmente-t-elle ?}
}
\end{minipage}

\clearpage

\canheadercorr{CAN 2}

\begin{minipage}[t]{0.493\textwidth}
\normalsize
\cantable{
\qrowcorr{1}{$72\div8$}{$9$}
\qrowcorr{2}{Écrire en chiffres : quinze unités et trois centièmes.}{$15{,}03$}
\qrowcorr{3}{$5-(-2)$}{$7$}
\qrowcorr{4}{Compléter : $47+\ldots=100$}{$53$}
\qrowcorr{5}{$10\,\%$ de 350}{$35$}
\qrowcorr{6}{La différence entre le double de 9 et le tiers de 12}{$14$}
\qrowcorr{7}{$0{,}8+1{,}35$}{$2{,}15$}
\qrowcorr{8}{Quel est le nombre qui, multiplié par 3, donne 7 ? Donner une fraction.}{$\dfrac73$}
\qrowcorr{9}{$3{,}6\div100$}{$0{,}036$}
\qrowcorr{10}{Dans la série 4 ; 6 ; 6 ; 7 ; 9, quel est l'effectif de 6 ?}{$2$}
\qrowcorr{11}{Le plus grand : $-3$ ou $-7$ ?}{$-3$}
\qrowcorr{12}{Quelle est la distance à zéro de $-6$ ?}{$6$}
\qrowcorr{13}{$\dfrac25$ de 35}{$14$}
\qrowcorr{14}{135 min = $\ldots$ h $\ldots$ min}{2 h 15 min}
\qrowcorr{15}{$(18-6)\div3$}{$4$}
}
\end{minipage}\hfill
\begin{minipage}[t]{0.493\textwidth}
\normalsize
\cantable{
\qrowcorr{16}{Le plus grand : $\dfrac34$ ou $\dfrac58$ ?}{$\dfrac34$}
\qrowcorr{17}{Aire d'un rectangle de 8 cm sur 3 cm}{24 cm$^2$}
\qrowcorr{18}{Moyenne de 7 ; 8 ; 9 ; 12}{$9$}
\qrowcorr{19}{Compléter : $1{,}2=\dfrac{\ldots}{10}$}{$12$}
\qrowcorr{20}{$1\,\mathrm{m}^2$ correspond à combien de carrés de $1\,\mathrm{cm}^2$ ?}{$10\,000$}
\qrowcorr{21}{Quel est le reste de la division de 29 par 4 ?}{$1$}
\qrowcorr{22}{$2{,}5\times4$}{$10$}
\qrowcorr{23}{$4\,500$ cm = $\ldots$ m}{45 m}
\qrowcorr{24}{4 cahiers coûtent 6 \texteuro. Combien coûtent 10 cahiers ?}{15 \texteuro}
\qrowcorr{25}{Dans un triangle, deux angles mesurent $90^\circ$ et $35^\circ$. Le troisième mesure :}{$55^\circ$}
\qrowcorr{26}{$11^2$}{$121$}
\qrowcorr{27}{Un carré de côté 8 cm peut-il avoir une aire de $70\,\mathrm{cm}^2$ ?}{non ($64\,\mathrm{cm}^2$)}
\qrowcorr{28}{Réduire : $x+x+x$}{$3x$}
\qrowcorr{29}{$5(x+2)$ est-il une somme ou un produit ?}{produit}
\qrowcorr{30}{La température passe de $-4\,^{\circ}\mathrm C$ à $3\,^{\circ}\mathrm C$. De combien augmente-t-elle ?}{$7\,^{\circ}\mathrm C$}
}
\end{minipage}
\end{document}
